\chapter{Marco normativo y ruta de cumplimiento}
\section{Jerarquía y forma de aplicar las normas}
Las leyes y reglamentos peruanos establecen obligaciones. El CNE fija requisitos de seguridad eléctrica; no sustituye los estudios hidráulicos, geotécnicos ni la selección de maquinaria. Las normas IEC, IEEE, ASME y de producto deben emplearse dentro de su alcance y con la edición contractual acordada. Una lista bibliográfica no demuestra cumplimiento.

Esta revisión vincula regla, criterio, evidencia y estado. Se distinguen la \textbf{corrección de una referencia}, la \textbf{especificación de una exigencia} y la \textbf{verificación de su satisfacción}. La matriz del Anexo A conserva abiertos los requisitos que dependen de campo, fabricante o autoridad.

\section{Código Nacional de Electricidad: alcance concreto}
\subsection{Suministro y utilización}
El CNE Suministro 2011 se aprobó mediante R.M. 214-2011-MEM/DM [N01]. Sus reglas 011.A-C incluyen generación y sistemas interactivos con suministro público, y establecen obligatoriedad en Perú. La regla 011.B remite las instalaciones del usuario en edificaciones al CNE Utilización. La central y su enlace requieren Suministro; las instalaciones interiores de servicios propios requieren compatibilizar Utilización y las reglas de estación aplicables.

El CNE Utilización fue aprobado por R.M. 037-2006-MEM/DM [N02]. Se leyó junto a la modificación R.M. 175-2008-MEM/DM [N03], especialmente propagación de fuego, diferenciales y puesta a tierra. Los factores de demanda de vivienda no se aplican automáticamente a bombas, gobernador, cargador o máquinas auxiliares.

\subsection{Reglas que cambian el diseño}
\begin{itemize}
\item \textbf{CNE-S 130.A, B y E:} además del regulador, considerar disparo independiente por sobrevelocidad cuando pueda resultar perjudicial, parada manual accesible y circuitos de parada protegidos.
\item \textbf{CNE-S 132:} el neutro de generador requiere resolver la impedancia y las necesidades particulares del sistema. No se adopta neutro sólido por analogía con una vivienda.
\item \textbf{CNE-S 036.A-B:} diseñar para paso, toque y potenciales transferidos. La regla 362 trata subidas de tuberías en poste; no sustenta un límite general de 5 ohm.
\item \textbf{CNE-S 110-124 y 152:} recinto, ventilación, humedad, iluminación, salidas, partes móviles, partes vivas y ubicación del transformador son verificaciones de diseño.
\item \textbf{CNE-U 030-004 y 050-104:} la sección del cable depende de instalación, temperatura, agrupamiento y régimen continuo; no basta dividir kVA entre tensión.
\item \textbf{CNE-U 050-102:} el 2.5\% individual y 4\% combinado corresponden a circuitos de utilización. La regulación del enlace de generación exige además requisitos de la distribuidora.
\end{itemize}

\subsection{Línea aérea y altitud}
Las distancias de la línea se evalúan en la Sección 23 (reglas 230, 232, 234 y tablas correspondientes). Las secciones 24, 25, 26 y 27 cubren grado de construcción, cargas, resistencia y aislamiento. Las secciones 21 y 22 no sustituyen esas tablas. Un poste de 12 m no acredita por sí mismo distancia al suelo ni resistencia.

La corrección por altitud exige leer el alcance de cada regla. Por ejemplo, 232.C.1.b actúa sobre la distancia adicional definida en 232.C.1.a para tensiones que sobrepasan 33 kV; no autoriza aumentar uniformemente toda distancia de una línea de 22.9 kV en 28\%. La coordinación de aislamiento externo de equipos se estudia con altitud, contaminación, humedad, impulso y especificación del fabricante; distancia de fuga y distancia en aire son verificaciones diferentes.

\section{Concesiones y obligaciones del titular}
El artículo 3 de la Ley de Concesiones Eléctricas utiliza el umbral \emph{mayor de 500 kW} para generación hidráulica. La inferencia para el objetivo de 150 kW es que no activa por sí solo ese supuesto de concesión definitiva. El artículo 7 conserva el cumplimiento técnico, ambiental, patrimonial y la obligación de informar inicio y características de operación [N04 y contraste N13]. La línea y las servidumbres tienen análisis propio; no se declara que todo el proyecto esté libre de autorizaciones.

El alumno es autor académico, no se le asigna automáticamente calidad de concesionario, propietario del agua o titular de obra. Antes de tramitar debe definirse titular, predios, modalidad de suministro, conexión y autoridad competente.

\section{Agua, caudal ecológico y obras}
La Ley 29338 regula derechos de uso. Sus artículos 44 y 47 sustentan el derecho y la licencia; el 53 exige disponibilidad en cantidad, calidad y oportunidad, preservación de caudales ecológicos y derechos de terceros, entre otros requisitos. El artículo 74 aborda faja marginal; el 81 la evaluación ambiental hídrica y el 104 las obras de infraestructura hidráulica [N05]. La acreditación de disponibilidad no equivale a licencia de uso ni a autorización para ejecutar obras.

El Reglamento de la Ley, artículos 153-155, reconoce permanencia y variación estacional del caudal ecológico [N06]. La R.J. 267-2019-ANA, artículo 6, vincula metodología a la categoría ambiental y establece un tratamiento para pequeños proyectos/no comprendidos en el listado SEIA coordinado con la AAA [N07]. Por tanto, no corresponde imponer ``20\% de Q95'' como regla universal.

\subsection{Secuencia de permisos y evidencia}
Se descargó el D.S. 023-2014-MINAGRI desde SPIJ [N19]. Sus artículos modificados 79, 81, 84 y 85 distinguen acreditación, autorización de obras y licencia. La acreditación no habilita a ejecutar obras ni usar agua. Las obras exigen antes el instrumento ambiental y la autorización sectorial cuando corresponde. La licencia requiere obras verificadas. No se atribuyen a la central procedimientos simplificados de otros tipos de proyecto.

\begin{longtable}{L{3.4cm}L{6.4cm}L{4.6cm}}
\toprule Etapa & Evidencia que debe prepararse & Condición del estudio\\\midrule
Disponibilidad & Balance por periodo, calidad, derechos, captación y restitución & No hay serie ni acreditación\\
Ambiente & Tipo de obra, evaluación del tramo derivado, caudal ecológico y certificación & Clasificación condicionada\\
Obras & Planos, predios, servidumbres, ambiente y requisito sectorial aplicable & Sin autorización ni planos ejecutivos\\
Licencia / operación & Verificación de obras y régimen de uso aprobado & Sin licencia\\
\bottomrule
\end{longtable}
El texto consultado es una modificación identificada, no consolidación de todas las reformas posteriores. La R.J. 007-2015-ANA y sus anexos deben acompañar el expediente administrativo; su publicación de aprobación no reemplaza los formatos completos. Antes de tramitar se confirma TUPA, autoridad y versión, sin deducir plazos vigentes de una copia histórica.

\section{Evaluación ambiental y territorio}
El D.S. 014-2019-EM, artículos 7-11 y Anexo 1, distingue estudios y clasificación [N08-N09]. En su texto base, una central de pasada sin estructura de regulación, de hasta 20 MW, ubicada en el supuesto B y con línea asociada de hasta 20 km, tiene clasificación anticipada DIA. Esta es una \textbf{ruta condicional}; no se ha demostrado que Limbani cumpla todas las condiciones B ni que la captación no genere embalse. No corresponde asignar DIA solo porque el objetivo es 150 kW.

Se verificará superposición con ANP y amortiguamiento, ACR, ecosistemas frágiles, derechos colectivos, población y patrimonio. Un barraje que genere embalse puede cambiar la fila aplicable del Anexo 1. El enlace y las obras auxiliares se evalúan como parte integral del proyecto, evitando fraccionar impactos.

La consulta pública de enero de 2026 es un proyecto normativo [N14]. Además se identificó el D.S. 011-2026-EM que modifica el artículo 48 para actividades en curso [N15]. Se incorpora esta actualización sin convertir un procedimiento de adecuación de instalaciones existentes en autorización para una central nueva. La consolidación normativa completa a la fecha de ejecución sigue pendiente.

\section{Obras civiles, seguridad y normas de equipo}
El RNE [N10] requiere seleccionar las versiones vigentes para cargas, sismo, suelos, concreto, acero, seguridad y electricidad. La geotecnia y el perfil definen cimentaciones, anclajes y deslizamiento; no se aprueban cuatro bloques por simple conteo del borrador.

La Ley 29783, su reglamento y el RESESATE aprobado por R.M. 111-2013-MEM/DM forman la base de seguridad laboral [N16-N17]. El plan debe desarrollar bloqueo de energías eléctrica e hidráulica, trabajo en altura, izaje, excavación, inundación, soldadura y rescate. Las referencias internacionales se especifican con acceso autorizado; la ficha IEC consultada ya publica edición 60034-1:2026 [N11]. No se traslada como obligación una fórmula de desclasificación sin texto ni curva de fabricante.
